\appendix
\section{재현에 필요한 입력 대응}\label{app:inputs}
PHYS16의 원 parent는 PHYS15이며 PR50의 확인된 head는
\begin{center}\small\texttt{b4b8019c6f149035c52784aced8e9cd7f3fbb23b}\end{center}
다. 입력의 권위는 파일 안에 기록된 exact endpoints와 binary64 bits다. 저장된 outward decimal은 읽기 위한 표시로 사용한다. 표 \ref{tab:inputmap}은 줄 번호가 바뀌어도 유지되는 JSON key 또는 source symbol로 대응한다.

\begin{table}[H]
\centering\small
\begin{tabularx}{\textwidth}{@{}p{44mm}X@{}}
\toprule
원 또는 신규 자료 & 읽거나 검증한 값\\\midrule
\path{UNIFORM_ARITHMETIC.json} & 같은 $C,B,q,S$; scalar $q$ 외에 원 componentwise $B$ 사용\\
\path{SOURCE_CENTRE.json} & \texttt{A\_Gg} 중심 Full residual Jacobian\\
\path{SOURCE_BOX.json} & \texttt{G\_gg}: Full residual Hessian; \texttt{q}: HH source jet\\
\path{SOURCE_*_COMPONENTS.json} & \texttt{nonphoto}: HH·외부 photo·팽창 추가 전 gas source jet\\
\path{C2_DOMAIN.json} & convex $X$, fraction 및 온도 domain; metadata와 실제 model threshold 구별\\
\path{ft03_interval_rhs} & photo-off 호출의 CI/RR/DR 값과 모든 thermal 수정 경로\\
\path{CORRELATED_DEFECT_RESULTS.json} & 공유 factor, 두 gas·일곱 photon residual, 35성분 secant, $d$와 readout\\
\path{ENDPOINT_REMAINDER_RESULTS.json} & $J,M,L_H,a$, $m=2$, sector 및 평균 adjoint 부호 포함\\
\path{WOLFRAM_XACT_EXECUTION.json} & 실패·보완 경로와 최종 여섯 tensor·네 scalar의 원 실행 응답\\
\bottomrule
\end{tabularx}
\caption{전체 상대 경로와 각 입력·산출물의 SHA256은 동봉 manifest, intake mapping 및 각 \texttt{RUN\_RECORD.json}에 보존된다. 원 PHYS15의 기존 science suite는 재실행하지 않았다.}
\label{tab:inputmap}
\end{table}

Physical coefficient의 상세 source 대응은 \texttt{SOURCE\_COOLING\_AUDIT\_KO.md}에 남겼다. 실제 threshold bits는
$\chi_H$: \texttt{402b3266051507a6},
$\chi_I$: \texttt{4038965f20504658},
$\chi_{II}$: \texttt{404b357928e0c9da}다.
각각 표시값은 $13.598434599702$, $24.587389011$, $54.41776$ eV다. $k_B$와 $\epsilon_{eV}$도 같은 captured binary64 leaf를 사용한다. 원 $\texttt{C2\_DOMAIN}$의 짧은 threshold 배열은 반올림 metadata여서 이 상수를 다시 정의하지 않는다.

논문의 $e$는 H당 전자 수이고 $n_e=ne$는 cm$^{-3}$다. $w$, $Q$, $U_\gamma$는 각각 eV/H, photons/H, eV/H로 정규화한다. $d$와 $\HF$는 같은 gas 단위를 가진다. Hessian $H^a_{jk}$는 $[g^a]/([g^j][g^k])$ 단위이며, $S$가 등장하는 노름에서 무차원화된다. $n_T,\overline z$는 온도 \emph{분자}를 위한 covector다. 이들을 $\Pi_F$로 나누기 전 값과 온도 기여를 혼동하지 않는다.

\section{문헌의 한 부호식에 대한 제한된 검산}\label{app:literaturesign}
확인한 Endtmayer--Langer--Wick의 RICAM 공개 preprint는 $\rho(\widetilde u)=-\mathcal A(\widetilde u)$로 정의한다. 식 (6)의 마지막 primal residual 항을 $-\rho(\widetilde u)(\widetilde z)$로 읽으면 다음 scalar 예가 일치하지 않는다.\cite{Endtmayer2019,lit16}
\begin{equation}
\mathcal A(u)(v)=uv,\quad J(u)=u,\quad
u=0,\quad z=1,\quad\widetilde u=1,\quad\widetilde z=1.
\end{equation}
여기서 모든 함수는 $C^\infty$이고 dual residual과 cubic remainder가 0이다. 정확한 goal 차이는 $-1$인데 해당 표시를 적용한 우변은 $+1$이다. $L(u,z)=J(u)-\mathcal A(u)(z)$에서 직접 유도하면 마지막 항은 $+\rho(\widetilde u)(\widetilde z)$다. 실제 scalar 검산 결과와 읽은 원문 위치는 문헌 패키지에 보존했다. 이 관찰은 접근한 공개 preprint의 식과 정의에 한정하며, 접근하지 못한 최종 journal 본문의 판정으로 확대하지 않는다. 본문의 endpoint·sector 정리는 해당 식과 독립적으로 증명했다.

\begin{thebibliography}{99}\small
\bibitem{phys14} WU088\_HH 연구 기록, \emph{PHYS14: geometric isotropic source-aligned 실제 endpoint 및 공통 arithmetic 포함}. PHYS15에 보존된 \path{inputs/phys14/} 원자료, root·source·arithmetic identity와 검토 기록.

\bibitem{phys15} WU088\_HH 연구 기록, \emph{PHYS15: 동일 COMMON의 정확한 Full--two-half residual과 공통 adjoint sector}. \path{HH_PHYS15_THEORY_20261012_v1}, PHYS15 보고서·contract·handoff. 확인된 PR50 head는 부록 \ref{app:inputs}에 기록했다.

\bibitem{app16} WU088\_HH PHYS16, \emph{실제 상관 endpoint residual 적용}. \path{evidence/intake/PHYS16_APPLICATION_DERIVATION_KO.md}, \path{checks/application/run01/results/CORRELATED_DEFECT_RESULTS.json}, \texttt{PRE\_RUN} 및 \texttt{RUN\_RECORD}. 2026-10-13 KST.

\bibitem{endpoint16} WU088\_HH PHYS16, \emph{Endpoint Jacobian과 이차 remainder에 의한 관측량 포함}. \path{theory/ENDPOINT_REMAINDER_THEORY_KO.md}; \path{checks/endpoint/ENDPOINT_HESSIAN_READOUT.json}; \path{checks/endpoint/application01/ENDPOINT_REMAINDER_RESULTS.json}. 2026-10-13 KST.

\bibitem{sourceaudit} WU088\_HH PHYS16, \emph{원 source cooling-cone 교차점검}. \path{evidence/intake/SOURCE_COOLING_AUDIT_KO.md}. 원 FT03 CI/RR/DR thermal symbol·상수·domain·export 순서와 source hash를 포함한 읽기 감사. 2026-10-13 KST.

\bibitem{review16} WU088\_HH PHYS16, \emph{독립 이론·산술 심사와 최종 채택}. \path{review/INDEPENDENT_REVIEW.json}, \path{review/INDEPENDENT_ARITHMETIC_RESULTS.json}, \path{review/independent_arithmetic_review.py}, \path{review/DECISION.json}. 18개 검산 묶음 통과, \texttt{PROMOTE\_SCOPED}, 2026-10-13 KST.

\bibitem{next16} WU088\_HH PHYS16, \emph{다음 매개변수 응답의 시작식과 입력 계약}. \path{theory/NEXT_PARAMETER_RESPONSE_PROPOSAL_KO.md}, \path{contracts/PARAMETER_RESPONSE_CONTRACT.json}, \path{NEXT_DAG.json}. 2026-10-13 KST.

\bibitem{Rump2010} S. M. Rump, ``Verification methods: Rigorous results using floating-point arithmetic,'' \emph{Acta Numerica} \textbf{19} (2010), 287--449. DOI: \href{https://doi.org/10.1017/S096249291000005X}{10.1017/S096249291000005X}. \href{https://www.tuhh.de/ti3/paper/rump/Ru10.pdf}{저자 공개 원문}; Theorems 10.2, 10.6, 10.8 및 Sections 11.3--11.6, 13.

\bibitem{Rannacher2003} R. Rannacher, ``Adaptive Finite Element Methods for Partial Differential Equations,'' arXiv:math/0305006 (2003), ICM 2002 lecture. \url{https://arxiv.org/pdf/math/0305006}. Propositions 1--2 및 residual--adjoint 전개.

\bibitem{Endtmayer2019} B. Endtmayer, U. Langer, T. Wick, ``Multigoal-oriented error estimates for non-linear problems,'' \emph{Journal of Numerical Mathematics} \textbf{27}(4) (2019), 215--236. DOI: \href{https://doi.org/10.1515/jnma-2018-0038}{10.1515/jnma-2018-0038}. 실제 확인 원문: \href{https://www.ricam.oeaw.ac.at/files/reports/18/rep18-08.pdf}{RICAM Report 2018-08}, Theorem 3.1과 Remarks 3.2--3.4. 부록 \ref{app:literaturesign}의 범위를 따른다.

\bibitem{Alekseev2005} A. K. Alekseev, I. M. Navon, ``Adjoint correction and bounding of error using Lagrange form of truncation term,'' \emph{Computers \& Mathematics with Applications} \textbf{50} (2005), 1311--1332. DOI: \href{https://doi.org/10.1016/j.camwa.2005.05.004}{10.1016/j.camwa.2005.05.004}. \href{https://people.sc.fsu.edu/~inavon/pubs/CAM2005.pdf}{저자 공개 원문}, Section 2.1.

\bibitem{xact} J. M. Martín-García 외, \emph{xAct / xTensor 공식 자료}. \url{https://xact.es/}; \url{https://xact.es/download.html}; \href{https://xact.es/Documentation/HTML/HTMLLinks/xTensor/ToCanonical.nb.html}{ToCanonical 문서}; \href{https://groups.google.com/g/xact/c/U6vGTjEZVDE}{xAct 개발자 지원의 내부 group backend 안내}. 2026-10-13 KST 확인.

\bibitem{xactrun} WU088\_HH PHYS16, \emph{Wolfram/xAct 실제 실행 기록}. \path{checks/wolfram/WOLFRAM_XACT_EXECUTION.json}, \path{checks/wolfram/endpoint_tensor_checks.wl}. 공식 archive SHA256:
\texttt{7a6c5f600868a3922668b020a15c0692f}\\
\texttt{76574ff2a559808c62d460cef1b07be}.

\bibitem{lit16} WU088\_HH PHYS16, \emph{문헌 조사와 endpoint bridge}. \path{literature/LITERATURE_REVIEW_KO.md}, \path{literature/ENDPOINT_BRIDGE_NOTE_KO.md}, 검색 provenance·manifest 및 \path{literature/PREPRINT_SIGN_CHECK.json}. Exa·SciSpace 검색과 1차 원문 검토, 2026-10-13 KST.
\end{thebibliography}
